\subsection{Pre-ordered monoids and groups}

Recall that if a relation \(x \preccurlyeq y\) between elements of a set E is reflexive and transitive, it is called a pre-order relation (Set Theory, III, p.~133). The relation « \(x \preccurlyeq y\) and \(y \preccurlyeq x\) » is an equivalence relation S on E, compatible with the relation \(x \preccurlyeq y\) ; on passing to the quotient, the relation \(\preccurlyeq\) induces an order relation on E/S, called the order relation associated to \(\preccurlyeq\) .

DEFINITION 2. --- A pre-order relation (written \(\preccurlyeq\) ) and a commutative monoid structure (written additively) on a set M are said to be compatible if they satisfy the following axiom:

(POM) For each \(z \in \mathbf{M}\) , the relation \(x \preccurlyeq y\) implies \(x + z \preccurlyeq y + z\) .

A set M equipped with a commutative monoid structure and a pre-order relation which are compatible is called a pre-ordered monoid.

Let M be a pre-ordered monoid, and S the equivalence relation \(\langle x \preccurlyeq y\) and \(y \preccurlyeq x \rangle\) . By virtue of (POM) the relation \(x \equiv x' \pmod{S}\) implies \(x + y \preccurlyeq x' + y\) and \(x' + y \preccurlyeq x + y\) for all \(y \in M\) , that is \(x + y \equiv x' + y \pmod{S}\) . In other words the equivalence relation S is compatible with addition in M (I, p.~11). Thus the quotient by S of the addition law on M, together with the ordering associated to \(\preccurlyeq\) , gives M/S the structure of an ordered monoid. In the case where M is a pre-ordered group, the group M/S is the quotient of M by the subgroup consisting of all elements x such that \(x \preccurlyeq 0\) and \(0 \preccurlyeq x\) .

